\subsection{Thảo luận về Các trường hợp Khó}
\label{sec:exp-scenarios}

\subsubsection*{Nhầm lẫn Web Attacks: XSS vs.\ Brute Force.}
Sự nhầm lẫn hệ thống giữa \texttt{Web Attack -- XSS} và \texttt{Web Attack -- Brute Force} là chế độ lỗi đáng kể nhất trong quy trình của chúng tôi xét về mặt vận hành. Cả hai loại tấn công đều nhắm vào lớp HTTP và tạo ra các Flow có thời lượng, số lượng Packet và phân phối IAT tương tự nhau. Ở cấp độ Flow, tín hiệu phân biệt duy nhất nằm ở các số liệu thống kê dẫn xuất từ Payload khá tinh vi (ví dụ: entropy của các kích thước phản hồi, phân phối độ dài URI) mà CICFlowMeter không nắm bắt được. Cả bốn mô hình trong thử nghiệm đồng thuận đều phân loại sai mẫu XSS thành Brute Force, xác nhận rằng sự nhầm lẫn này không dành riêng cho mô hình nào mà vốn có trong Feature Space. Một biện pháp giảm thiểu thực tế sẽ là gộp hai loại phụ của Web Attacks thành một lớp \texttt{Web Attack} duy nhất để tạo cảnh báo, chấp nhận một nhãn thô hơn để đổi lấy khả năng phát hiện đáng tin cậy.

\subsubsection*{Rare\_Attack: Nhóm không đồng nhất.}
Nhóm \texttt{Rare\_Attack} gộp ba loại tấn công có dấu hiệu mạng khác nhau về chất: \texttt{Infiltration} (Beacon đánh cắp dữ liệu khối lượng thấp), \texttt{SQL Injection} (các Flow nặng về HTTP POST) và \texttt{Heartbleed} (dị thường ở lớp bản ghi TLS cụ thể). Việc gộp chúng thành một nhãn duy nhất buộc bộ phân loại phải học một ranh giới quyết định duy nhất cho ba phân phối riêng biệt, điều này chắc chắn làm tăng sự nhầm lẫn với cả lớp \texttt{BENIGN} và lớp Web Attacks. Thử nghiệm đồng thuận cho thấy hai mô hình dự đoán đúng \texttt{Rare\_Attack} trong khi hai mô hình còn lại dự đoán là \texttt{Web Attack -- XSS} và \texttt{BENIGN}, phản ánh sự không đồng nhất này. Với nhiều dữ liệu hơn (hoặc các đặc trưng chuyên dụng ở cấp độ Payload), các lớp này sẽ xứng đáng có các nhãn riêng và đánh giá mô hình cá nhân.

\subsubsection*{Phát hiện Bot.}
Lớp \texttt{Bot} (1.948 mẫu trong ngữ liệu gốc) đạt được $F_1 = 0.86$ với LightGBM sau khi áp dụng SMOTE. Lưu lượng bot trong CIC-IDS2017 được tạo bởi bản phát lại PCAP của botnet Ares, tạo ra các Beacon C2 dựa trên HTTP đặc trưng với thời gian xuất hiện định kỳ. Mức độ recall tương đối cao (0.80) cho thấy mô hình nắm bắt được tính đều đặn về mặt thời gian, nhưng tỷ lệ khoảng 80\% dữ liệu tổng hợp đã gây ra một số nhiễu nhãn làm giảm độ chính xác (precision). Việc triển khai IDS với ngưỡng quyết định thấp hơn đối với lớp Bot (đánh đổi với việc FPR trên lưu lượng BENIGN tăng nhẹ) sẽ làm tăng recall về phía giá trị vận hành mong muốn là $> 0.95$.
